\documentclass[UTF8]{ctexart}
\usepackage{amsmath}
\usepackage{tocloft}
\usepackage[bigfiles]{pdfbase}
\usepackage{amsthm,amssymb}
\usepackage{booktabs}
\usepackage{graphicx}
\usepackage[hidelinks]{hyperref}
\usepackage{geometry}
\usepackage{float}
\usepackage{multirow}
\usepackage{indentfirst}
\usepackage{diagbox}
\usepackage{listings}
\usepackage{xcolor}
\definecolor{codegreen}{rgb}{0,0.6,0}
\definecolor{codegray}{rgb}{0.5,0.5,0.5}
\definecolor{codepurple}{rgb}{0.58,0,0.82}
\definecolor{backcolour}{rgb}{0.95,0.95,0.92}

\lstdefinestyle{mystyle}{
    backgroundcolor=\color{backcolour},   
    commentstyle=\color{codegreen},
    keywordstyle=\color{magenta},
    numberstyle=\tiny\color{codegray},
    stringstyle=\color{codepurple},
    basicstyle=\ttfamily\footnotesize,
    breakatwhitespace=false,         
    breaklines=true,                 
    captionpos=b,                    
    keepspaces=true,                 
    numbers=left,                    
    numbersep=5pt,                  
    showspaces=false,                
    showstringspaces=false,
    showtabs=false,                  
    tabsize=2
}
\lstset{style=mystyle}
\newcommand\question[2]{\noindent\fbox{\parbox[t]{\linewidth}{\hspace{0pt}{\textbf{#1\quad}#2}}}}


\usepackage{graphicx}
\usepackage{setspace}
\renewcommand{\cftsecleader}{\cftdotfill{\cftdotsep}}
\geometry{a4paper,left=3cm,right=3cm,top=2cm,bottom=2cm} 


\CTEXsetup[format={\Large \bfseries}]{section}
\title{第一次上机作业：YOLOv5}
\author{陈天乐\quad 无25\quad [student ID removed]\\ \href{mailto:ctl22@mails.tsinghua.edu.cn}{ctl22@mails.tsinghua.edu.cn}}
\date{\today}
\begin{document}
\maketitle
\tableofcontents
\newpage

\section{背景}
YOLOv5 is the world’s most loved vision AI。本次上机作业要求使用官方提供的 YOLOv5 COCO 预训练权重（不再训练模型）和已经给出的 367 张测试图片和对应的标注（COCO 目标类别）完成推理任务。

\section{作业任务}
\subsection{使用 YOLOv5s 完成推理，计算检测性能}
将官网上的yolov5模型下载至本地后，编写了一个coco\_test.yaml文件，用于定义测试验证的图片数据集，同时由于annotations.json中的id和category都是不连续的，因此编写了一个Label.py用于读出annotations.json中的
id和对应的category，写入coco\_test.yaml中，代码在报告结尾给出。在终端调用val.py程序对给出的测试图片进行验证:
\begin{lstlisting}[frame=shadowbox]
    python val.py --weight yolov5s.pt --data ./data/coco_test.yaml --task val --img 640 --batch-size 1
\end{lstlisting}

得到mAP@50和mAP@50:5:95的结果如下图所示：

\begin{figure}[H]
    \centering
    \includegraphics[width=\textwidth]{v5s.png}
    \caption{使用 YOLOv5s 完成推理的检测性能}
\end{figure}

\subsection{可视化预测的结果，并分析典型失败案例}
在终端调用detect.py对给出的图片数据集进行推理，得到带有预测框的输出图片，并将label保存到一个txt文件中：
\begin{lstlisting}[frame=shadowbox]
    python detect.py --source ImageSet --weights yolov5s.pt --img 640 --save-txt 
\end{lstlisting}

然后编写了一个Draw\_GT.py文件，在带有预测框的输出图片中将annotations.json的GT框画上，代码在报告结尾给出。这里给出几个典型失败案例和其分析原因（白色框为GT框，其余颜色是推理框）：

\begin{figure}[H]
    \centering
    \includegraphics[width=\textwidth]{C:/Users/skyhappy/Desktop/学校的破事/作业/数字图像处理/DIP_assignment/DIP_HW/visual_results/026cde425ea95a4d.jpg}
    \caption{失败案例1}
\end{figure}

模型对这个图片发生了虚警，可能是由于标注的class中没有老虎这一类别，因此模型将该图片检测为了bear。

\begin{figure}[H]
    \centering
    \includegraphics[width=\textwidth]{C:/Users/skyhappy/Desktop/学校的破事/作业/数字图像处理/DIP_assignment/DIP_HW/visual_results/6a638756c8168376.jpg}
    \caption{失败案例2}
\end{figure}

在这个图片中，模型正确识别了图中的人，但是在中间这个女人身上还有衣服和背包两种物品未能识别出来，这是由于YOLO的算法导致的，对于很多重叠在一起的物体时，YOLO只能对
一整个物体进行识别，而不能对这个物体重叠部分的小物体进行识别。

\begin{figure}[H]
    \centering
    \includegraphics[width=\textwidth]{C:/Users/skyhappy/Desktop/学校的破事/作业/数字图像处理/DIP_assignment/DIP_HW/visual_results/58295753041c640b.jpg}
    \caption{失败案例3}
\end{figure}

在这张图片中，物体分布的非常紧凑，同时还有一些非常细小的物体被标注为GT。该模型基本能够把大部分物体识别出来，但是对于一些非常小型、细节的物体
未能识别，这同样也是由于YOLO算法的原因。

\subsection{调整可能调整的配置}
对这批图片使用了不带TTA的模型（n/m/l/x/s6)和带TTA的模型（n/m/l/x)，运行结果如下图：
\begin{figure}[H]
    \centering
    \includegraphics[width=\textwidth]{v5n.png}
    \caption{YOLOv5n模型的检测性能}
\end{figure}
\begin{figure}[H]
    \centering
    \includegraphics[width=\textwidth]{v5m.png}
    \caption{YOLOv5m模型的检测性能}
\end{figure}
\begin{figure}[H]
    \centering
    \includegraphics[width=\textwidth]{v5l.png}
    \caption{YOLOv5l模型的检测性能}
\end{figure}
\begin{figure}[H]
    \centering
    \includegraphics[width=\textwidth]{v5x.png}
    \caption{YOLOv5x模型的检测性能}
\end{figure}
\begin{figure}[H]
    \centering
    \includegraphics[width=\textwidth]{v5s6.png}
    \caption{YOLOv5s6模型的检测性能}
\end{figure}
\begin{figure}[H]
    \centering
    \includegraphics[width=\textwidth]{v5n_aug.png}
    \caption{带TTA的YOLOv5n模型的检测性能}
\end{figure}
\begin{figure}[H]
    \centering
    \includegraphics[width=\textwidth]{v5m_aug.png}
    \caption{带TTA的YOLOv5m模型的检测性能}
\end{figure}
\begin{figure}[H]
    \centering
    \includegraphics[width=\textwidth]{v5l_aug.png}
    \caption{带TTA的YOLOv5l模型的检测性能}
\end{figure}
\begin{figure}[H]
    \centering
    \includegraphics[width=\textwidth]{v5x_aug.png}
    \caption{带TTA的YOLOv5x模型的检测性能}
\end{figure}

将mAP@50和mAP@50:5:95的数据汇总为下表：
\begin{table}[!ht]
    \centering
    \renewcommand{\arraystretch}{1.5}
    \caption{不带TTA的模型检测精度}
    \begin{tabular}{|c|c|c|}
    \hline
        模型 & mAP@50 & mAP@50:5:95 \\ \hline
        YOLOv5s & 0.379 & 0.272 \\ \hline
        YOLOv5n & 0.287 & 0.200 \\ \hline
        YOLOv5m & 0.406 & 0.316 \\ \hline
        YOLOv5l & 0.442 & 0.358 \\ \hline
        YOLOv5x & 0.462 & 0.374 \\ \hline
    \end{tabular}

\end{table}
\begin{table}[!ht]
    \centering
    \renewcommand{\arraystretch}{1.5}
    \caption{带TTA的模型检测精度}
    \begin{tabular}{|c|c|c|}
    \hline
        模型 & mAP@50 & mAP@50:5:95 \\ \hline
        YOLOv5s & 0.384 & 0.297 \\ \hline
        YOLOv5x & 0.423 & 0.341 \\ \hline
    \end{tabular}
\end{table}

观察统计数据后可以发现，更换不同模型能够显著影响检测精度，因为不同模型的参数量大小不同,但对于该图片测试集而言，并不是
参数量最大的x模型表现最好，反而是m模型表现较好，这可能是模型过拟合导致的。同时，带有TTA的模型只有n型有较明显的提高，剩余的甚至会有下降，但是基本不变。调整模型的conf\_thres后，对模型的检测精度
影响也不是非常明显。

\subsection{提出改进建议}
根据yolo之后的版本以及yolov5出现的问题，我认为其可以从以下几个部分提高：
\subsubsection{主干网络（Backbone）优化}
参考YOLOv8的改进，其引入了轻量化视觉变换器（RepViT），结合了CNN的局部感知能力和ViT的全局建模能力，能在减少参数量的同时提升检测精度，尤其适合处理复杂背景和小目标场景，这很大程度上
能缓解YOLOv5对小物体难以识别的困难。同时，参考YOLOv6多尺度特征融合的特点，主干网络可以采用跨阶段局部网络（CSPNet）优化特征提取效率，结合快速空间金字塔池化（SPPF）模块，既保留多尺度信息，又降低计算复杂度。



\subsubsection{注意力机制增强}
参考YOLOv8的改进，在主干网络的关键位置插入轻量级Transformer块，结合局部卷积与全局自注意力机制，可以提升复杂场景下的特征表达能力
，这样遇到复杂背景图片时检测性能可以大大提高。同时，Transformer的一大特点就是注意力机制，在Backbone中加入位置坐标注意力，通过动态调整特征权重，增强关键区域的响应，同样可以提升
微小目标的检测精度。

\subsubsection{锚框匹配策略优化}
参考 YOLOv7 的动态正样本选择策略改进，可以优化Anchors的匹配策略，调整锚框尺寸以适应不同数据集，例如为微小目标设计更密集的锚框，采用动态匹配规则，增加正样本数量以加快收敛速率。


\subsubsection{轻量化与速度优化}
观察到YOLOv5不同网络的推理速度在CPU上运行还需要几百毫秒，对于实时应用而言延时可能还是会有点大。因此可以对网络进行轻量化处理，比如
进行量化、剪枝操作，还可以在量化时采用动态量化技术（如INT8量化），在保证精度的前提下压缩模型体积，提升推理速度，适用于边端部署和实时系统。


\newpage
\section{源代码}
\begin{lstlisting}[frame=shadowbox,caption = Label.py,language=python]
import json
import time
import os
from tqdm import tqdm
from ultralytics import YOLO


# COCO 官方80类的原始ID（非连续列表）
COCO_OFFICIAL_IDS = [
    1, 2, 3, 4, 5, 6, 7, 8, 9, 10, 11, 13, 14, 15, 16, 17, 18, 19, 20,
    21, 22, 23, 24, 25, 27, 28, 31, 32, 33, 34, 35, 36, 37, 38, 39, 40,
    41, 42, 43, 44, 46, 47, 48, 49, 50, 51, 52, 53, 54, 55, 56, 57, 58,
    59, 60, 61, 62, 63, 64, 65, 67, 70, 72, 73, 74, 75, 76, 77, 78, 79,
    80, 81, 82, 84, 85, 86, 87, 88, 89, 90
]

# 配置参数（根据实际情况修改）
annotations_json_path = "C:/Users/skyhappy/Desktop/学校的破事/作业/数字图像处理/DIP_assignment/annotations.json" # COCO标注文件路径
output_path = "C:/Users/skyhappy/Desktop/DIP_HW" # 输出目录
image_path = "C:/Users/skyhappy/Desktop/DIP_HW/images" # 图片目录
val_path = os.path.join(output_path, 'image.txt')


# 创建映射字典：{原始COCO_ID : 连续YOLO_ID (0-79)}
COCO_TO_YOLO_ID = {coco_id: yolo_id for yolo_id, coco_id in enumerate(COCO_OFFICIAL_IDS)}

def coco2yolo(annotation_json_path, output_path, image_patj):
    # 创建输出目录
    labels_dir = os.path.join(output_path, 'Label')
    os.makedirs(labels_dir, exist_ok=True)


# 加载COCO标注
    with open(annotations_json_path, 'r') as f:
    coco_data = json.load(f)
    image_entries = coco_data.get("images", [])

    images_info = {img['id']: img for img in coco_data['images']}
    
    # 按图像分组标注
    annotations = {}
    for ann in coco_data['annotations']:
        img_id = ann['image_id']
        if img_id not in annotations:
            annotations[img_id] = []
        annotations[img_id].append(ann)

    for img_id, img_info in tqdm(images_info.items(), desc="Processing Images"):
        # 获取图片信息
        img_name = img_info['file_name']
        img_width = img_info['width']
        img_height = img_info['height']
    
        # 验证图片存在
        img_path = os.path.join(image_path, img_name)
        if not os.path.exists(img_path):
            print(f"Warning: Image {img_name} not found, skipping")
            continue
        
        with open(val_path, 'w') as f:
            for entry in image_entries:
            name = entry.get("file_name")
            f.write(f"C:/Users/skyhappy/Desktop/DIP_HW/images/{name}\n")
        f.close

        # 创建对应的标签文件
        label_name = os.path.splitext(img_name)[0] + '.txt'
        label_path = os.path.join(labels_dir, label_name)
    
        # 写入YOLO格式标注
        with open(label_path, 'w') as f:
            if img_id in annotations:
                for ann in annotations[img_id]:
                    # 检查是否为COCO 80类
                    coco_id = ann['category_id']
                    if coco_id not in COCO_TO_YOLO_ID:
                        continue # 跳过非COCO80类的标注

                    # 获取连续YOLO ID
                    class_id = COCO_TO_YOLO_ID[coco_id]
        
                    # 转换边界框格式,归一化
                    x, y, w, h = ann['bbox']
                    x_center = (x + w / 2) / img_width
                    y_center = (y + h / 2) / img_height
                    w_norm = w / img_width
                    h_norm = h / img_height
        
                    # 写入：class_id x_center y_center width height
                    f.write(f"{class_id} {x_center:.6f} {y_center:.6f} {w_norm:.6f} {h_norm:.6f}\n")

if __name__ == "__main__":
    # 执行转换
    coco2yolo(annotations_json_path, output_path, image_path)
    print(f"Conversion complete! Labels saved to: {os.path.join(output_path, 'Label')}")
\end{lstlisting}

\begin{lstlisting}[frame=shadowbox,caption = Draw\_GT.py,language=python]
import json
import os
import shutil
from pathlib import Path
import numpy as np
import cv2
from tqdm import tqdm
import argparse
import yaml


def parse_args():
    parser = argparse.ArgumentParser(description='YOLOv5检测结果可视化工具')
    parser.add_argument('--img-dir', type=str, default='C:/Users/skyhappy/Desktop/DIP_HW/exp3', help='预测输出图片目录')
    parser.add_argument('--label-dir', type=str, default='C:/Users/skyhappy/Desktop/DIP_HW/Label', help='真实标签目录')
    parser.add_argument('--output-dir', type=str, default='C:/Users/skyhappy/Desktop/DIP_HW/visual_results', help='可视化结果输出目录')
    parser.add_argument('--yaml-path', type=str, default='C:/Users/skyhappy/Desktop/DIP_HW/coco_test.yaml', help='数据集配置文件路径')
    parser.add_argument('--gt-color', type=tuple, default=(255, 255, 255), help='真实框颜色（BGR）')
    return parser.parse_args()


def load_class_names(yaml_path):
    """从YAML文件加载类别名称"""
    with open(yaml_path, 'r', encoding='utf-8') as f:
        data = yaml.safe_load(f)
    return data['names']

def draw_boxes(image, label_path, color, class_names, is_gt=True):
    h, w = image.shape[:2]
    with open(label_path, 'r', encoding='utf-8') as f:
        for line in f:
            parts = line.strip().split()
            if not parts:
                continue
            
            try:
                if is_gt:
                    cls_id, xc, yc, bw, bh = map(float, parts)
                    conf = None
                else:
                    cls_id, xc, yc, bw, bh, conf = map(float, parts)
            except Exception as e:
                print(f"错误：标签文件 {label_path} 中存在格式错误的行：{line.strip()}，错误信息：{e}")
                continue

            # 坐标转换
            x1 = int((xc - bw/2) * w)
            y1 = int((yc - bh/2) * h)
            x2 = int((xc + bw/2) * w)
            y2 = int((yc + bh/2) * h)
            # 绘制边界框
            cv2.rectangle(image, (x1, y1), (x2, y2), color, 2)

            if not is_gt:
                cls_id = int(cls_id)
                if cls_id >= len(class_names):
                    print(f"警告：预测类别ID {cls_id} 超出类别列表范围（共{len(class_names)}类）")
                    continue
                label = f"{class_names[cls_id]} {conf:.2f}"
                (tw, th), _ = cv2.getTextSize(label, cv2.FONT_HERSHEY_SIMPLEX, 0.6, 2)
                cv2.rectangle(image, (x1, y1-20), (x1+tw, y1), color, -1)
                cv2.putText(image, label, (x1, y1-5), cv2.FONT_HERSHEY_SIMPLEX, 0.6, (255,255,255), 2)
    
    return image

def main(args):
    os.makedirs(args.output_dir, exist_ok=True)
    class_names = load_class_names(args.yaml_path)
    processed = 0
    for img_name in os.listdir(args.img_dir):
        if not img_name.lower().endswith(('.jpg', '.jpeg', '.png')):
            continue

        base_name = os.path.splitext(img_name)[0]
        img_path = os.path.join(args.img_dir, img_name)
        gt_label_path = os.path.join(args.label_dir, f"{base_name}.txt")
        pred_label_path = os.path.join(args.pred_dir, f"{base_name}.txt")

        print(f"\n处理图片：{img_name}")
        print(f"真实标签路径：{gt_label_path}")
        print(f"预测标签路径：{pred_label_path}")

        image = cv2.imread(img_path)
        if image is None:
            print(f"错误：无法读取图片 {img_path}，请检查文件格式或路径。")
            continue

        image = draw_boxes(image, gt_label_path, args.gt_color, class_names, is_gt=True)

        output_path = os.path.join(args.output_dir, img_name)
        cv2.imwrite(output_path, image)
        processed += 1
        print(f"成功保存：{output_path}")


if __name__ == "__main__":
    args = parse_args()
    main(args)
\end{lstlisting}

\begin{lstlisting}[frame=shadowbox,caption = coco\_test.yaml,language=python]
# Ultralytics 🚀 AGPL-3.0 License - https://ultralytics.com/license
# parent
# ├── yolov5
# └── datasets
#     └── coco_test

# Train/val/test sets as 1) dir: path/to/imgs, 2) file: path/to/imgs.txt, or 3) list: [path/to/imgs1, path/to/imgs2, ..]
path: C:/Users/skyhappy/Desktop/DIP_HW # dataset root dir
train: image.txt # train images (relative to 'path') 
val:  image.txt # val images (relative to 'path') 
test: image.txt # test images (optional)
labels_dir: Label

# Classes
names:
    0: person
    1: bicycle
    2: car
    3: motorcycle
    4: airplane
    5: bus
    6: train
    7: truck
    8: boat
    9: traffic light
    10: fire hydrant
    11: stop sign
    12: parking meter
    13: bench
    14: bird
    15: cat
    16: dog
    17: horse
    18: sheep
    19: cow
    20: elephant
    21: bear
    22: zebra
    23: giraffe
    24: backpack
    25: umbrella
    26: handbag
    27: tie
    28: suitcase
    29: frisbee
    30: skis
    31: snowboard
    32: sports ball
    33: kite
    34: baseball bat
    35: baseball glove
    36: skateboard
    37: surfboard
    38: tennis racket
    39: bottle
    40: wine glass
    41: cup
    42: fork
    43: knife
    44: spoon
    45: bowl
    46: banana
    47: apple
    48: sandwich
    49: orange
    50: broccoli
    51: carrot
    52: hot dog
    53: pizza
    54: donut
    55: cake
    56: chair
    57: couch
    58: potted plant
    59: bed
    60: dining table
    61: toilet
    62: tv
    63: laptop
    64: mouse
    65: remote
    66: keyboard
    67: cell phone
    68: microwave
    69: oven
    70: toaster
    71: sink
    72: refrigerator
    73: book
    74: clock
    75: vase
    76: scissors
    77: teddy bear
    78: hair drier
    79: toothbrush
\end{lstlisting}


\end{document}
